% =========================================================
% Irrational Arithmetic
% =========================================================

\section{Irrational Arithmetic}
\label{sec:irrational-arithmetic}

These results record how irrationality behaves under arithmetic with rational
numbers, and why the irrational numbers fail to form a field.

\begin{tcolorbox}[colback=thmbox, colframe=thmborder, arc=2pt,
  left=6pt, right=6pt, top=4pt, bottom=4pt,
  title={\small\textbf{Theorem (Rational Plus Irrational Is Irrational)}},
  fonttitle=\small\bfseries]
\begin{theorem}[Rational Plus Irrational Is Irrational]
\label{thm:rational-plus-irrational-is-irrational}
If $q\in\mathbb{Q}$ and $\alpha\in\mathbb{R}\setminus\mathbb{Q}$, then
\[
q+\alpha\in\mathbb{R}\setminus\mathbb{Q}.
\]
\end{theorem}
\end{tcolorbox}
\begin{remark*}[Dependencies]
Closure of $\mathbb{Q}$ under subtraction.
\end{remark*}
\begin{remark*}[Proof idea]
If $q+\alpha$ were rational, subtracting the rational number $q$ would force
$\alpha$ to be rational.
\end{remark*}
\begin{remark*}[Proof]
\hyperref[prf:rational-plus-irrational-is-irrational]{Go to proof.}
\end{remark*}

\begin{tcolorbox}[colback=thmbox, colframe=thmborder, arc=2pt,
  left=6pt, right=6pt, top=4pt, bottom=4pt,
  title={\small\textbf{Theorem (Rational Minus Irrational Is Irrational)}},
  fonttitle=\small\bfseries]
\begin{theorem}[Rational Minus Irrational Is Irrational]
\label{thm:rational-minus-irrational-is-irrational}
If $q\in\mathbb{Q}$ and $\alpha\in\mathbb{R}\setminus\mathbb{Q}$, then
\[
q-\alpha\in\mathbb{R}\setminus\mathbb{Q}.
\]
\end{theorem}
\end{tcolorbox}
\begin{remark*}[Dependencies]
Closure of $\mathbb{Q}$ under subtraction.
\end{remark*}
\begin{remark*}[Proof idea]
If $q-\alpha$ were rational, then subtracting it from $q$ would force
$\alpha$ to be rational.
\end{remark*}
\begin{remark*}[Proof]
\hyperref[prf:rational-minus-irrational-is-irrational]{Go to proof.}
\end{remark*}

\begin{tcolorbox}[colback=thmbox, colframe=thmborder, arc=2pt,
  left=6pt, right=6pt, top=4pt, bottom=4pt,
  title={\small\textbf{Theorem (Nonzero Rational Times Irrational Is Irrational)}},
  fonttitle=\small\bfseries]
\begin{theorem}[Nonzero Rational Times Irrational Is Irrational]
\label{thm:nonzero-rational-times-irrational-is-irrational}
If $q\in\mathbb{Q}\setminus\{0\}$ and $\alpha\in\mathbb{R}\setminus\mathbb{Q}$,
then
\[
q\alpha\in\mathbb{R}\setminus\mathbb{Q}.
\]
\end{theorem}
\end{tcolorbox}
\begin{remark*}[Dependencies]
Nonzero rational numbers have rational reciprocals.
\end{remark*}
\begin{remark*}[Proof idea]
If $q\alpha$ were rational, divide by the nonzero rational $q$ to conclude that
$\alpha$ is rational.
\end{remark*}
\begin{remark*}[Proof]
\hyperref[prf:nonzero-rational-times-irrational-is-irrational]{Go to proof.}
\end{remark*}

\begin{tcolorbox}[colback=thmbox, colframe=thmborder, arc=2pt,
  left=6pt, right=6pt, top=4pt, bottom=4pt,
  title={\small\textbf{Theorem (Quotient by Nonzero Rational Preserves Irrationality)}},
  fonttitle=\small\bfseries]
\begin{theorem}[Quotient by Nonzero Rational Preserves Irrationality]
\label{thm:quotient-by-nonzero-rational-preserves-irrationality}
If $\alpha\in\mathbb{R}\setminus\mathbb{Q}$ and
$q\in\mathbb{Q}\setminus\{0\}$, then
\[
\frac{\alpha}{q}\in\mathbb{R}\setminus\mathbb{Q}.
\]
\end{theorem}
\end{tcolorbox}
\begin{remark*}[Dependencies]
Nonzero rational numbers have rational reciprocals.
\end{remark*}
\begin{remark*}[Proof idea]
If $\alpha/q$ were rational, then multiplying by $q$ would force $\alpha$ to
be rational.
\end{remark*}
\begin{remark*}[Proof]
\hyperref[prf:quotient-by-nonzero-rational-preserves-irrationality]{Go to proof.}
\end{remark*}

\begin{tcolorbox}[colback=corbox, colframe=corborder, arc=2pt,
  left=6pt, right=6pt, top=4pt, bottom=4pt,
  title={\small\textbf{Corollary (The Irrationals Are Not Closed Under Addition)}},
  fonttitle=\small\bfseries]
\begin{corollary}[The Irrationals Are Not Closed Under Addition]
\label{cor:irrationals-not-closed-under-addition}
The set $\mathbb{R}\setminus\mathbb{Q}$ is not closed under addition.
\end{corollary}
\end{tcolorbox}
\begin{remark*}[Dependencies]
\hyperref[thm:rational-minus-irrational-is-irrational]{Theorem~\ref*{thm:rational-minus-irrational-is-irrational}}.
\end{remark*}
\begin{remark*}[Proof idea]
Choose irrational numbers whose sum is rational, for example
$\sqrt{2}$ and $2-\sqrt{2}$.
\end{remark*}
\begin{remark*}[Proof]
\hyperref[prf:irrationals-not-closed-under-addition]{Go to proof.}
\end{remark*}

\begin{tcolorbox}[colback=corbox, colframe=corborder, arc=2pt,
  left=6pt, right=6pt, top=4pt, bottom=4pt,
  title={\small\textbf{Corollary (The Irrationals Are Not Closed Under Multiplication)}},
  fonttitle=\small\bfseries]
\begin{corollary}[The Irrationals Are Not Closed Under Multiplication]
\label{cor:irrationals-not-closed-under-multiplication}
The set $\mathbb{R}\setminus\mathbb{Q}$ is not closed under multiplication.
\end{corollary}
\end{tcolorbox}
\begin{remark*}[Dependencies]
The irrationality of $\sqrt{2}$.
\end{remark*}
\begin{remark*}[Proof idea]
Use $\sqrt{2}\cdot\sqrt{2}=2$.
\end{remark*}
\begin{remark*}[Proof]
\hyperref[prf:irrationals-not-closed-under-multiplication]{Go to proof.}
\end{remark*}

\begin{tcolorbox}[colback=corbox, colframe=corborder, arc=2pt,
  left=6pt, right=6pt, top=4pt, bottom=4pt,
  title={\small\textbf{Corollary (The Irrationals Are Not a Field)}},
  fonttitle=\small\bfseries]
\begin{corollary}[The Irrationals Are Not a Field]
\label{cor:irrationals-are-not-a-field}
The set $\mathbb{R}\setminus\mathbb{Q}$ does not form a field under the usual
operations.
\end{corollary}
\end{tcolorbox}
\begin{remark*}[Dependencies]
\hyperref[cor:irrationals-not-closed-under-addition]{Corollary~\ref*{cor:irrationals-not-closed-under-addition}},
\hyperref[cor:irrationals-not-closed-under-multiplication]{Corollary~\ref*{cor:irrationals-not-closed-under-multiplication}}.
\end{remark*}
\begin{remark*}[Proof idea]
Failure of closure under addition or multiplication already prevents a set from
being a field.
\end{remark*}
\begin{remark*}[Proof]
\hyperref[prf:irrationals-are-not-a-field]{Go to proof.}
\end{remark*}

\begin{tcolorbox}[colback=corbox, colframe=corborder, arc=2pt,
  left=6pt, right=6pt, top=4pt, bottom=4pt,
  title={\small\textbf{Corollary (The Irrationals Are Not an Ordered Field)}},
  fonttitle=\small\bfseries]
\begin{corollary}[The Irrationals Are Not an Ordered Field]
\label{cor:irrationals-are-not-an-ordered-field}
The set $\mathbb{R}\setminus\mathbb{Q}$ does not form an ordered field under the
usual operations and order.
\end{corollary}
\end{tcolorbox}
\begin{remark*}[Dependencies]
\hyperref[cor:irrationals-are-not-a-field]{Corollary~\ref*{cor:irrationals-are-not-a-field}}.
\end{remark*}
\begin{remark*}[Proof idea]
An ordered field is, in particular, a field.
\end{remark*}
\begin{remark*}[Proof]
\hyperref[prf:irrationals-are-not-an-ordered-field]{Go to proof.}
\end{remark*}
